\SourceLeaf{268}
\begin{list}{}{\setlength{\leftmargin}{1.5em}\setlength{\itemindent}{-1.5em}\setlength{\itemsep}{0pt}\setlength{\parsep}{0pt}}
\item \(e\) is the eccentricity of the ellipse that generates the ellipsoidal surface by revolving about its minor axis.
\item \(t\) and \(90^0-w\) are the longitude and latitude of an arbitrary point on this surface. Thus \(w\) is the angle between the normal to the surface at this point and the minor axis.
\item \(T\) and \(90^0-U\) are the longitude and latitude of the corresponding point on the spherical surface.
\item \(i\) is the imaginary unit \(\BareRoot{-1}\).
\item \(f\) is the symbol for a function that can be chosen arbitrarily.
\end{list}
\noindent The logarithms are always natural.\par
\noindent We shall denote the scale factor by \(m\). It is defined so that each line element on the ellipsoidal surface stands to the corresponding element on the spherical surface as 1 to \(m\). The ratio has a definite value at each point of either surface and can vary from point to point.

The simplest solution is obtained by taking the arbitrary function to be equal to its argument, setting
\[
f\upsilon=\upsilon
\]
This is indeed the most suitable mapping if the \textit{entire} surface of the ellipsoid is to be mapped onto the spherical surface. In geodetic calculations, where only a relatively small part of the Earth’s surface is ever considered, it is much more advantageous, as already noted there, to add a constant imaginary part to the function, setting
\[
f\upsilon=\upsilon-i\log k
\]
The radius of the sphere and the constant \(k\) can then be chosen so that the scale factor \(m\) \(=1\) on the middle parallel and differs scarcely at all from 1 within a few degrees to the north and south. The convenience of using the mapping in geodetic calculations depends chiefly on how little this factor varies within the area mapped. The deviation from 1 is of the second order with respect to
